\section{Introduction}
\label{intro}

CERN runs a large computing infrastructure. Since 2012 the IT department has built up a configuration management (CM) system around Puppet and OpenStack, together with a number of in-house services, to manage its computing services~\cite{bellinfra2012}. Each of these tools does its own job well. What is missing is anything that ties them together, and the cost shows up in traceability, reproducibility, and maintainability.

Procedures for managing servers do exist. Keeping the service documentation in step with the infrastructure is the harder problem, particularly in an environment with frequent personnel rotation. When the person who set a service up is no longer around, recovering or reconfiguring it is slow. Because IaC tools are declarative, the desired state of the infrastructure is the documentation, which removes much of that dependency on institutional memory. There is a second gap as well: Puppet and tools like it manage the configuration of a server, not the server itself~\cite{loope2011managing}. When the machine underneath is degraded or gone, CM has nothing to converge.

Over the past few years IaC has become the standard model for consolidating how infrastructure is developed, deployed, and automated~\cite{morris2025infrastructure}. Terraform and its community-driven fork, OpenTofu, are the most widely used implementations, and both are a plausible basis for unifying provisioning and configuration in a reproducible and scalable way~\cite{holkuziev2025infrastructure}.

This work sets out to design and implement an IaC-based framework that integrates with the existing infrastructure management ecosystem rather than replacing it, and that gives operators a single automated and traceable way to manage resources spread across cloud regions. If the integration works, the payoff is less operational complexity and maintenance effort, better automation and reproducibility, and a far easier path through Business Continuity / Disaster Recovery (BC/DR) and cross-environment replication. We see it as a first step towards a more unified infrastructure management strategy at CERN.
